\section{Evaluations des indicateurs de risque}

Dans cette partie, nous nous intéresserons dans un premier temps à la théorie entourant les indicateurs de risque, appelés \textbf{Grecques}, puis nous utiliserons le modèle binomial pour essayer de les approximer. Nous nous limiterons ici aux grecques principales : $\Delta, \Gamma, \nu, \Theta$. Il en existe de nombreuses autres, notamment $\rho, \lambda, \epsilon$, mais qui sont moins pertinentes en raison des hypothèses présupposées du modèle binomial (taux sans risque constant, pas de dividende, etc.).

\subsection{Delta $(\Delta)$}

\subsection{Gamma $(\Gamma)$}

\subsection{Vega $(\nu)$}

\subsection{Thêta $(\Theta)$}